\section{Introduction}\label{sec:introduction}

Reliable daily streamflow simulation remains a central problem in
surface-water hydrology, particularly in basins where observations are
sparse, hydroclimatic variability is strong and operational decisions
must still be supported by defensible discharge estimates
\citep{razavi2013streamflow,guo2021regionalization}. These conditions
are characteristic of many high-elevation Andean catchments. In the
Peruvian Altiplano, precipitation is strongly seasonal, monitoring
networks are sparse relative to the spatial heterogeneity of rainfall,
and discharge dynamics are shaped by the alternation between sharp
wet-season runoff response and sustained dry-season low flows
\citep{condom2004transient,zubieta2018hydrological,lima2025modeling}.
Such settings expose a persistent limitation of classical
rainfall--runoff modelling: even when a conceptual model captures the
main water-balance controls, simplified storage structures, uncertain
parameters and imperfect meteorological inputs can lead to structural
errors that become particularly visible under low-flow, flashy or
climatically contrasting conditions
\citep{van2013influence,nauditt2017conceptual,ragettli2014evaluation}.

Machine-learning (ML) models offer a complementary route because they
can learn nonlinear rainfall--runoff relations directly from data. Large
sample studies have shown that recurrent and other data-driven models
can match or outperform lumped conceptual models in daily streamflow
prediction, including in regional and basin-scale benchmarks
\citep{kratzert2019towards,lees2021benchmarking,arsenault2022continuous}.
However, predictive skill alone does not eliminate hydrological risk.
Purely data-driven models may be sensitive to distributional shifts,
training-period representativeness, extrapolation to extremes and the
absence of explicit physical constraints
\citep{frame2021deep,zhong2023developing}. In data-scarce basins, their
performance also depends strongly on whether antecedent hydrological
state information is available, because rainfall inputs alone may be
insufficient to represent catchment memory
\citep{adnan2021comparison,okkan2021embedding}. The practical question
is therefore not whether process-based or ML models are universally
superior, but how physical structure and empirical flexibility can be
combined without obscuring the source of any improvement.

Hybrid physical--ML modelling addresses this question by using process
simulations, physically derived states or synthetic physically
consistent samples as inputs to an empirical emulator. This strategy has
improved runoff simulation in several configurations, including
conceptual-model correction, nested hybrid rainfall--runoff modelling,
physics-guided deep learning and process-driven deep learning
\citep{mekonnen2015hybrid,mohammadi2022ihacres,okkan2021embedding,yang2020physical,xie2021physics,li2024process,zhu2025hybrid}.
Most of these studies, however, pass forward a deterministic physical
trajectory. That design preserves part of the process information but
collapses the conditional distribution of the simulated response into a
single value. For stochastic rainfall--runoff transformation, this is a
strong simplification: uncertainty in precipitation, model structure and
hydrological memory propagates nonlinearly into discharge and can alter
both peak timing and magnitude
\citep{gabellani2007propagation,hong2006uncertainty}. It also limits the
ability of the hybrid model to exploit distributional descriptors such
as ensemble quantiles, which are closely related to hydrological
signatures and to probabilistic forecasting diagnostics
\citep{guo2021regionalization,papacharalampous2022review}.

A recent stochastic hybridization study by
\citet{houenafa2025hybridization} showed that mean and quantile
descriptors of a L\'{e}vy-driven stochastic HyMoLAP discharge ensemble
can improve ML-based rainfall--runoff prediction. This result is
important because it moves physical--ML hybridization from deterministic
state correction toward distribution-aware physical descriptors. Yet it
also leaves three issues unresolved. First, the gain from the hybrid
pipeline may come from several components---hydrological memory,
stochastic perturbations, the ML correction itself or the use of
quantile rather than mean ensemble descriptors---and these effects are
not identifiable from a single best-model comparison. Second, a
stochastic ensemble that improves point prediction is not necessarily a
calibrated predictive interval; coverage, sharpness and interval scores
must be evaluated separately. Third, reproducibility claims require that
reported model labels match the backend that actually produced the
results, especially when deep-learning architectures have software
fallbacks or implementation contingencies.

This study addresses these gaps through StoHyMoLAP, an audited
stochastic physical--ML framework for daily streamflow simulation in the
high-Andean Ramis basin, the principal tributary of Lake Titicaca. The
framework couples a lumped RAMIS rainfall--runoff core, an optional
linear baseflow reservoir, L\'{e}vy-type stochastic perturbations,
Monte Carlo physical ensembles and ML emulators trained either from
observed forcings or from physical-ensemble descriptors. The purpose is
not to introduce another best-performing hybrid model, but to decompose
which components are responsible for performance gains. Nine controlled
experiments (E0--E8; Table~\ref{tab:experiment-matrix}) are therefore
evaluated over a common validation window using a single workflow
(Fig.~\ref{fig:method-workflow}). The ablation design isolates the
incremental roles of baseflow memory, stochastic forcing, ML correction
and descriptor choice.

The manuscript makes four specific contributions. First, it transfers
stochastic physical--ML quantile hybridization to a high-Andean,
strongly seasonal and data-scarce Altiplano basin, providing a contrast
to previous applications outside this hydroclimatic setting. Second, it
shows whether ensemble mean or ensemble quantiles carry the more useful
physical information for ML correction. Third, it evaluates the same
stochastic ensemble both as a source of predictive features and as a
candidate uncertainty band through PICP, PINAW and Winkler diagnostics.
Fourth, it embeds temporal-leakage and backend audits into the evidence
chain, following the broader call for transparent and reproducible
hydrological modelling \citep{hall2021hydrologist}. A key audit finding
is stated explicitly: the experiments labelled as GRU in the current
run were executed with a fallback multilayer perceptron on flattened
sequences, not with a true recurrent backend. They are therefore
interpreted as GRU-labelled sequential hybrids under fallback execution,
not as evidence of recurrent-network superiority
\citep{kratzert2018rainfall}.

%% =====================================================================
%% 2. MATERIALS AND METHODS
%% =====================================================================
